\subsection*{\S. 14.\quad 可展曲面。}
在空间曲线理论中，我们只需考察一种情形：被包络面构成一个单参数族，且全都是平面。在这种情形下，包络面称为\emph{可展曲面}；这个名称的理由马上说明。

平面族中每一个平面的特征线显然是一条直线，而所有特征直线都是回归棱的切线；因此，可展曲面是一条曲线的所有切线的轨迹，这条曲线就是它的回归棱。不难看出，动平面（被包络面）与回归棱的密切平面重合。事实上，对这条曲线采用通常的记号，其密切平面方程为
\[
(X-x)\lambda+(Y-y)\mu+(Z-z)\nu=0
\]
动密切平面上的特征线，由这个方程与它对 \(s\) 微分所得的方程联立而得，即
\[
(X-x)\xi+(Y-y)\eta+(Z-z)\zeta=0;
\]
而这两个平面的交线恰好就是切线。

不过应当注意，回归棱可能退化为一个点；这时，按照这个点位于有限距离处还是无穷远处，可展曲面分别成为锥面或柱面。

每一个被包络平面都沿整条特征直线（母线）与可展曲面相切，因此可展曲面的切平面构成一个单参数族；而对任何其他曲面，切平面都构成一个双参数族\srcfn{1}{根据对偶原理，每一个曲面都对应于一个曲面；只有可展曲面例外，它对应于一条曲线。}。

“可展”这一名称源于如下事实：若设曲面可以弯曲但不能伸长，则可以把它展开到平面上而不发生撕裂或重叠。反过来，具有这一性质的任何曲面必定是可展曲面，这一点以后将予以证明。

与每一条空间曲线有关，需要考察三个可展曲面；它们分别由主三面形的三个面包络而成。密切平面的包络，就是给定曲线全部切线的轨迹，如上所述。法平面的包络称为曲线 C 的\emph{极可展面}；垂直于主法线的各个平面的包络，则称为\emph{矫直可展面}。
